\section{Meta-closure against same-domain incompleteness}
\label{sec:meta-closure}

Section~\ref{sec:physics-closure} closed the physical side of the capstone. The next and last pre-synthesis burden is to show that the fixed domain cannot be reopened from above by code, truth, infinitary rule-power, or richer same-scope extension. The claim is exact. This section does not deny classical G\"odel incompleteness in its ordinary arithmetized setting. It proves only that a same-domain G\"odel architecture threatening governance-level closure is not instantiable inside the unique admissible interior already fixed in Section~\ref{sec:standing-regime}.\footnote{This is the precise capstone use of \noninstgodeldoc{} \cite{MaleyNonInst}: same-domain non-instantiability, not denial of classical metamathematics.}

The local theorem support for this section is now straightforward: the dedicated same-domain non-instantiability manuscript supplies the meta-closure package, while Sections~\ref{sec:modeling-necessity} and \ref{sec:standing-regime} provide the already-fixed trajectory and standing architecture on which the present route-exhaustion argument depends.\cite{MaleyNonInst}

\subsection{Inferential life and closure relevance}

\begin{definition}[Inferential life of the fixed domain]
The \emph{inferential life} of the fixed domain is the totality of what may be admitted, rejected, committed, reused, or licensed as same-scope continuation on that domain.
\end{definition}

\begin{theorem}[Inferential-life exhaustion]
\label{thm:inferential-life-exhaustion}
On the fixed scope already determined by Sections~\ref{sec:modeling-necessity} and \ref{sec:standing-regime}, inferential life is exhausted by exactly two structures:
\begin{enumerate}
  \item the standing classification of evaluated acts on the fixed domain; and
  \item the terminal standing behavior of legal same-scope trajectories.
\end{enumerate}
More precisely, if two same-domain descriptions agree on the standing classification of every evaluated act and on the terminal standing behavior of every legal same-scope trajectory, then they agree on what may be admitted, rejected, committed, reused, or licensed as same-scope continuation.
\end{theorem}

\begin{proof}
At the act level, Section~\ref{sec:standing-regime} already fixed a unique same-scope gate, standing emergence as reuse-stable admissibility, standing-quotient uniqueness, and the no-second-same-scope-equivalence theorem. So once standing classification of evaluated acts is fixed, no additional same-domain act-level standing-bearing distinction remains.

At the path level, Section~\ref{sec:modeling-necessity} fixed coherence as standing preservation along legal same-scope trajectories and proved that incoherent prefixes remain incoherent under extension. Standing failure is therefore irrecoverable within the retained lineage. So once terminal standing behavior of legal same-scope trajectories is fixed, what may be licensed as same-scope continuation or commitment is fixed as well.

Since inferential life consists exactly of what may be admitted, rejected, committed, reused, or licensed as same-scope continuation, agreement on those two structures exhausts it.
\end{proof}

\begin{theorem}[Closure-relevance transmission]
\label{thm:closure-relevance-transmission}
Let $\Theta$ be any same-domain distinction on the fixed scope. If $\Theta$ changes inferential life, then $\Theta$ changes at least one of the following:
\begin{enumerate}
  \item the standing classification of some evaluated act; or
  \item the terminal standing behavior, and hence the licensed continuation structure, of some legal same-scope trajectory.
\end{enumerate}
Consequently, every same-domain incompleteness objection can threaten closure only if it survives into the standing quotient or the coherent trajectory structure of the fixed domain.
\end{theorem}

\begin{proof}
Assume that $\Theta$ changes neither the standing classification of any evaluated act nor the terminal standing behavior of any legal same-scope trajectory. Then, by Theorem~\ref{thm:inferential-life-exhaustion}, two descriptions differing only by $\Theta$ agree on the entire inferential life of the fixed domain. That contradicts the hypothesis that $\Theta$ changes inferential life. So at least one of the two listed kinds of change must occur.

If the change is act-level, then by standing-quotient uniqueness it survives into the standing quotient: there is no second same-scope equivalence relation that can absorb it while leaving admissible standing unchanged. If the change is trajectory-level, then by the finite fixed-scope trajectory results it survives into coherent trajectory structure itself, because no later same-scope continuation can repair an incoherent prefix. Hence every same-domain closure threat must survive into one of those two structures.
\end{proof}

\begin{corollary}[External undecidability inertness]
An undecidable or unprovable sentence that never re-enters the fixed scope so as to change standing classification or coherent same-scope trajectory structure is closure-inert on that scope.
\end{corollary}

\begin{proof}
If neither kind of change occurs, then Theorem~\ref{thm:closure-relevance-transmission} shows that the distinction changes no inferential-life fact and therefore threatens no same-domain closure.
\end{proof}

\begin{theorem}[Inferential necessity of closure distinctions]
\label{thm:inferential-necessity-closure-distinctions}
Let $F$ be any formal system associated with the fixed domain, and let $\Theta$ be a sentence, code-construction, or metamathematical distinction formulated in $F$. Suppose the claim is that the undecidability, unprovability, or fixed-point behavior of $\Theta$ undercuts closure of reasoning on the fixed domain in the ordinary inferential sense, namely that it changes what may be admitted, rejected, committed, or licensed as same-domain continuation when $\Theta$ is used in reasoning on that domain. Then $\Theta$ must induce at least one of the following:
\begin{enumerate}
  \item a difference in the standing classification of some evaluated act; or
  \item a difference in the terminal standing of at least one legal same-domain trajectory, equivalently a divergence in what may be licensed as standing-bearing continuation or commitment on the fixed domain.
\end{enumerate}
Consequently, every ordinary G\"odel-style phenomenon can threaten closure on the fixed domain only if it survives into the standing quotient or coherent trajectory structure of that domain.
\end{theorem}

\begin{proof}
By hypothesis, the presence or absence of $\Theta$ changes the inferential life of the fixed domain. Therefore Theorem~\ref{thm:closure-relevance-transmission} yields either a standing-classification difference or a trajectory-level terminal-standing difference.

If the difference is act-level, then by standing-quotient uniqueness and the collapse of second same-scope standing-bearing equivalence relations it survives into the standing quotient itself. If the difference is trajectory-level, then by the finite fixed-scope trajectory theorems it survives into coherent trajectory structure, because terminal standing certifies coherence, incoherent prefixes remain incoherent under extension, and standing failure cannot be repaired downstream within the same lineage. Therefore every ordinary G\"odel-style phenomenon that genuinely threatens closure must survive into the standing quotient or coherent trajectory structure of the fixed domain.
\end{proof}

\begin{corollary}[Trajectory-neutral closure relevance is act-level]
If an alleged distinction on the fixed domain alters no legal same-domain trajectory---equivalently, changes neither any terminal standing outcome nor any licensed same-domain continuation---then any closure relevance it has must appear as a difference in standing classification of some evaluated act.
\end{corollary}

\begin{proof}
Apply Theorem~\ref{thm:closure-relevance-transmission}. If the trajectory branch is absent by hypothesis, only the act-level branch remains.
\end{proof}

\begin{remark}
\label{rem:godel-governance-relevance}
Why G\"odel must become governance-relevant to matter. Classical incompleteness by itself is not enough to threaten closure in the present sense. To matter here, an undecidable or unprovable sentence must change the inferential life of the fixed domain; and by Theorems~\ref{thm:closure-relevance-transmission} and \ref{thm:inferential-necessity-closure-distinctions}, any such change must survive into standing classification or coherent trajectory structure. The bridge from standard metamathematical talk to governance-level closure is therefore theorematic rather than interpretive.
\end{remark}

\subsection{Governance-relevant code predicates are hidden gate work}

\begin{definition}[Same-domain faithful coding layer]
A \emph{same-domain faithful coding layer} for a class $A$ of syntax-bearing acts on the fixed scope consists of a coding map
\[
\ulcorner\cdot\urcorner : A \to C
\]
into a code-object space $C$ such that predicates on $C$ may be read back, without scope reset or domain change, as standing-bearing constraints on the original acts or on their licensed same-scope continuations.
\end{definition}

\begin{definition}[Readback and governance-relevant code predicate]
If $P\subseteq C$ is a predicate on code-objects, its same-domain readback is the induced predicate
\[
\widehat P(a) :\iff P(\ulcorner a\urcorner)
\qquad (a\in A).
\]
The code predicate $P$ is \emph{governance-relevant} when this readback changes inferential life on the fixed domain.
\end{definition}

\begin{lemma}[Hidden gate-work lemma]
\label{lem:hidden-gate-work}
If a same-domain code predicate is governance-relevant, then its readback performs hidden gate work on the original fixed domain: it either changes which acts are admitted or standing-positive, or it refines the positive side or negative side of the gate into a further substantive same-scope classification affecting licensed continuation.
\end{lemma}

\begin{proof}
If the readback is governance-relevant, then by definition it changes inferential life. Theorem~\ref{thm:closure-relevance-transmission} therefore forces the change to appear either in standing classification or in coherent trajectory structure. If it changes standing classification directly, it changes which acts are admitted or standing-positive. If it leaves admission unchanged but changes what same-scope continuation is licensed, then it refines one side of the gate into a further substantive same-scope classification. No third same-domain option remains.
\end{proof}

\begin{theorem}[Code-predicate role exhaustion]
\label{thm:code-role-exhaustion}
Let $P\subseteq C$ be a same-domain code predicate with readback $\widehat P$. Then exactly one of the following roles is instantiated on the fixed domain:
\begin{enumerate}
  \item inert bookkeeping: $\widehat P$ changes no standing classification and no licensed same-scope continuation;
  \item extensional gate renaming: $\widehat P$ coincides extensionally with the unique same-scope admissibility gate and therefore adds no independent governance content;
  \item same-side refinement: $\widehat P$ leaves gate verdicts unchanged but purports to refine the admitted side or rejected side into a further substantive same-scope classification affecting licensed continuation;
  \item independent standing-bearing authorizer: $\widehat P$ purports to authorize admissibility or standing-bearing continuation independently of the unique same-scope gate and is not extensionally identical with it.
\end{enumerate}
No fifth same-domain governance role exists.
\end{theorem}

\begin{proof}
If $\widehat P$ is closure-inert, then role (1) holds by definition. Assume instead that $\widehat P$ is governance-relevant. By Lemma~\ref{lem:hidden-gate-work}, it either changes who is admitted or standing-positive, or it leaves admission fixed while changing licensed same-scope continuation. In the first case, either the induced classification agrees extensionally with the unique gate, giving role (2), or it does not, giving role (4). In the second case it refines one side of the gate without changing admission, giving role (3). These cases are mutually exclusive and exhaustive.
\end{proof}

\begin{theorem}[No independent governance-relevant code predicate]
\label{thm:no-independent-code-predicate}
No same-domain code predicate can serve as an independent governance-relevant authorizer on the original fixed scope. Every same-domain code predicate is therefore either inert bookkeeping or extensional renaming of the unique admissibility gate.
\end{theorem}

\begin{proof}
By Theorem~\ref{thm:code-role-exhaustion}, every same-domain code predicate must instantiate exactly one of the four roles listed there. Role (3) is impossible because Section~\ref{sec:standing-regime} already eliminated substantive same-side refinements: a positive-side refinement would amount to a second same-scope gate, while a negative-side refinement would open a same-scope promotion channel, both of which are excluded by the standing theorem family. Role (4) is impossible because the no-carriers theorem of Section~\ref{sec:standing-regime} already excludes any independent same-scope authorizer not extensionally identical with the unique gate. So only roles (1) and (2) remain.
\end{proof}

\begin{corollary}[One numbering is not enough]
A single fixed coding map does not evade Theorem~\ref{thm:no-independent-code-predicate}. One numbering is harmless unless predicates on its codes re-enter the fixed scope as governance-relevant constraints; and once they do, code-role exhaustion applies.
\end{corollary}

\begin{proof}
The danger is governance-relevant readback, not multiplicity of coding as such. A single numbering that never changes standing or trajectories is inert bookkeeping; a single numbering that does change them is already covered by Theorems~\ref{thm:code-role-exhaustion} and \ref{thm:no-independent-code-predicate}.
\end{proof}

\subsection{Same-domain diagonalization and the gate-collapsed fixed point}

\begin{definition}[Same-domain negative fixed-point construction]
A \emph{same-domain negative fixed-point construction} on the fixed scope consists of:
\begin{enumerate}
  \item a same-domain faithful coding layer on a class $A$ of syntax-bearing acts;
  \item a governance-relevant code predicate $P\subseteq C$ with readback $\widehat P$;
  \item an admissible same-domain substitution/readback operation producing an act $g\in A$ whose standing-bearing content is lawfully equivalent to a negative fixed-point formula about its own code, typically of the form $\neg P(\ulcorner g\urcorner)$.
\end{enumerate}
\end{definition}

\begin{theorem}[Necessary package for a same-domain G\"odel threat]
\label{thm:same-domain-godel-prerequisites}
If a G\"odel-style undecidability, unprovability, or fixed-point phenomenon is claimed to threaten governance-level closure on the fixed scope from within the same domain, then any such threat must realize all of the following simultaneously:
\begin{enumerate}
  \item a same-domain faithful coding layer;
  \item governance-relevant readback of at least one code-based distinction to the original fixed domain;
  \item a standing-bearing code-predicate role that survives the code-role exhaustion of this section;
  \item an admissible same-domain fixed-point production; and
  \item survival of the resulting distinction into the standing quotient or coherent trajectory structure of the fixed domain.
\end{enumerate}
\end{theorem}

\begin{proof}
Without same-domain faithful coding, the phenomenon remains external bookkeeping. Without governance-relevant readback, it changes no standing fact and no same-scope continuation fact, so it is closure-inert by the previous subsection. Without a standing-bearing code-predicate role, there is no way for the code-level distinction to matter to the fixed domain at all. Without an admissible same-domain fixed-point production, there is no G\"odel-style self-reference on the original scope. And if the resulting distinction does not survive into the standing quotient or coherent trajectory structure, then by Theorem~\ref{thm:closure-relevance-transmission} it does not threaten closure. So every same-domain G\"odel threat requires all five items.
\end{proof}

\begin{theorem}[Non-existence of admissible same-domain diagonalization]
\label{thm:no-same-domain-diagonalization}
No same-domain diagonal apparatus producing a governance-relevant negative fixed point is instantiable on the fixed scope.
\end{theorem}

\begin{proof}
Assume, for contradiction, that a same-domain diagonal apparatus exists. By Theorem~\ref{thm:same-domain-godel-prerequisites}, it must realize a same-domain faithful coding layer, governance-relevant readback, a standing-bearing code-predicate role, an admissible negative fixed-point production, and survival of the resulting distinction into the standing quotient or coherent trajectory structure. By Theorem~\ref{thm:no-independent-code-predicate}, the only remaining governance-relevant code-level case after role exhaustion is extensional gate renaming. So the standing-bearing content of the fixed point $g$ is lawfully equivalent to ``$g$ is not admitted'' or ``$g$ is not standing-positive.''

If $g$ is admitted or standing-positive, then the same evaluated act is both admitted and, by its own standing-bearing content, not admitted. That contradicts the unique same-scope gate and the bivalent standing theorem family of Section~\ref{sec:standing-regime}.

If $g$ is not admitted or not standing-positive, then it carries no standing-bearing force on the fixed scope. For it nevertheless to threaten closure, its failed status would have to be reversed or made governance-relevant again either by same-act repair / generation or by later same-scope continuation within the same lineage. The first route is blocked by the no-repair and no-generator theorems of Section~\ref{sec:standing-regime}; the second is blocked by the no-deferred-coherence and irrecoverability results of Section~\ref{sec:modeling-necessity}. In either case the purported fixed point becomes closure-inert. That contradicts item (5) of Theorem~\ref{thm:same-domain-godel-prerequisites}. Therefore no admissible same-domain diagonalization exists.
\end{proof}

\begin{corollary}[No governance-relevant negative fixed point]
There is no governance-relevant fixed-domain act $g$ whose standing-bearing content is lawfully equivalent to $\neg\mathrm{Prov}(\ulcorner g\urcorner)$ for any same-domain provability predicate $\mathrm{Prov}$.
\end{corollary}

\begin{proof}
Such an act would be a same-domain negative fixed point and therefore part of an admissible same-domain diagonal apparatus. Theorem~\ref{thm:no-same-domain-diagonalization} eliminates that possibility.
\end{proof}


\subsubsection*{Inline proof anchor VI: why same-domain diagonalization fails}

The full same-domain obstruction is proved in \noninstgodeldoc{} \cite{MaleyNonInst}. Once the fixed-domain standing package is already in place, the diagonal failure argument is finite.
\begin{enumerate}
  \item Any same-domain G\"odel threat must realize the full prerequisite package: a faithful coding layer, governance-relevant readback, a standing-bearing code-predicate role, an admissible fixed-point production, and survival of the resulting distinction into standing classification or coherent trajectory structure.
  \item Code-role exhaustion leaves only four roles for the governing code-predicate: inert bookkeeping, extensional gate renaming, same-side refinement, and independent authorizer. The last two are impossible on the fixed domain, so the only governance-relevant surviving case is extensional gate renaming.
  \item The putative negative fixed point is therefore gate-collapsed. Its standing-bearing content is lawfully equivalent not to an independent provability split but simply to ``not admitted'' or ``not standing-positive.''
  \item If the fixed-point act $g$ is admitted, then the same evaluated act is both admitted and, by its own standing-bearing content, not admitted. That contradicts bivalence together with unique gate classification.
  \item If $g$ is not admitted, then the construction threatens closure only if its failed status can later regain same-domain governance force. Same-act repair and same-source generation are already blocked on the fixed domain, and same-lineage downstream restoration is blocked by irrecoverability of standing failure. So the non-admitted branch is closure-inert.
\end{enumerate}
Every same-domain negative fixed point therefore collapses into one of two dead ends: contradiction on the admitted branch or inertness on the non-admitted branch. That is why Theorem~\ref{thm:no-same-domain-diagonalization} is stronger than the slogan ``G\"odel does not apply here'': it shows that the relevant same-domain apparatus never reaches a stable governance-relevant fixed point inside the unique admissible interior.

\subsection{Semantic, infinitary, and extension rescue routes}

\begin{theorem}[No semantic or meta-foundational standing import]
\label{thm:no-semantic-standing-import}
Truth predicates, satisfaction predicates, intended-model semantics, global closure predicates, reflection-on-truth, and similar same-domain semantic or meta-foundational resources cannot function as independent same-domain standing-bearing authorizers on the fixed scope. Any such appeal is therefore either closure-inert bookkeeping, extensional gate renaming, explicit scope change, or inadmissible same-domain authority import.
\end{theorem}

\begin{proof}
If such a semantic or meta-foundational appeal changes no standing classification and no coherent same-scope trajectory fact, then by Theorem~\ref{thm:closure-relevance-transmission} it is closure-inert bookkeeping. If it agrees extensionally with the unique same-scope gate, it is only gate renaming. If it purports to alter admissibility, standing, or licensed continuation while remaining on the old scope, then it either functions as an independent same-scope authorizer or attempts to preserve a second closure-relevant equivalence relation alongside admissible standing classification. Section~\ref{sec:standing-regime} excludes the first; the same section's no-second-equivalence theorem excludes the second. So the only remaining non-inert use is explicit scope change or inadmissible same-domain authority import.
\end{proof}

\begin{definition}[Finite same-domain witness]
A \emph{finite same-domain witness} for a purported same-domain distinction $\Theta$ is either an evaluated act whose standing classification differs when $\Theta$ is imposed, or a finite legal same-scope trajectory whose coherence, terminal standing, or licensed continuation status differs when $\Theta$ is imposed.
\end{definition}

\begin{theorem}[Infinitary admissibility collapse]
\label{thm:infinitary-admissibility-collapse}
Let $\Theta$ be a same-domain standing-bearing distinction attributed to reflection, infinitary rule-power, or transfinite closure. If $\Theta$ is governance-relevant on the fixed scope, then $\Theta$ has a finite same-domain witness. Equivalently, every admissible same-domain standing-bearing use of infinitary rule-power collapses to finite witness at the level of evaluated acts or finite legal trajectories; otherwise the appeal is closure-inert or an inadmissible same-domain authorizer.
\end{theorem}

\begin{proof}
If $\Theta$ is governance-relevant, then by Theorem~\ref{thm:closure-relevance-transmission} it must survive either into standing classification or into coherent trajectory structure. In the first case, some evaluated act changes standing status and is therefore a finite same-domain witness. In the second case, some legal same-scope trajectory changes coherence, terminal standing, or licensed continuation. But Section~\ref{sec:modeling-necessity} fixes the trajectory theory only for finite nonempty same-scope trajectories, so this is again a finite same-domain witness. Conversely, if no such finite witness exists, then no standing classification changes and no finite legal same-scope trajectory changes; the alleged infinitary distinction is therefore closure-inert or else an attempt to smuggle in a new same-domain authorizer.
\end{proof}

\begin{corollary}[No primitive infinitary same-domain standing acts]
There is no primitive same-domain act-token whose alleged standing-bearing force depends on an infinitary verification burden not reduced to a finite same-domain witness already visible on the fixed scope.
\end{corollary}

\begin{proof}
If the alleged infinitary burden changes no standing or trajectory fact beyond the ordinary gate verdict on the act-token, then it is bookkeeping only. If it is claimed to matter to governance-level closure, then by Theorem~\ref{thm:infinitary-admissibility-collapse} there must already be a finite same-domain witness. If one denies such a witness while still claiming same-domain standing-bearing force, then the token can do that work only by semantic or meta-foundational import, which Theorem~\ref{thm:no-semantic-standing-import} excludes.
\end{proof}

\begin{theorem}[Same-domain extension trichotomy]
\label{thm:same-domain-extension-trichotomy}
Let $E$ be any attempt to rescue a same-domain incompleteness objection by adjoining a richer theory, additional primitives, stronger rules, or auxiliary structure. Then exactly one of the following holds:
\begin{enumerate}
  \item $E$ changes no standing classification and no coherent trajectory fact on the old scope, in which case it is definitional or bookkeeping only;
  \item $E$ operates only after explicit reset, fresh gating, or domain change, in which case it is an explicit scope change rather than a same-domain rescue; or
  \item $E$ purports to change standing classification or coherent trajectory structure on the old scope, in which case it is inadmissible same-domain authority import.
\end{enumerate}
In particular, there is no fourth same-domain extension route to a governance-relevant incompleteness threat.
\end{theorem}

\begin{proof}
If the extension changes no standing classification and no coherent same-scope trajectory fact on the old scope, then by Theorem~\ref{thm:inferential-life-exhaustion} it changes no inferential-life fact at all and is therefore bookkeeping only. If it works only after explicit reset, fresh gating, or domain change, then by definition it is a scope change rather than a same-domain rescue. Assume finally that it remains on the old scope while changing standing classification or coherent trajectory structure there. Then Theorem~\ref{thm:closure-relevance-transmission} forces the difference to survive into the standing quotient or coherent trajectory structure. But Section~\ref{sec:standing-regime} already fixed the standing quotient as unique and eliminated any second same-scope standing-bearing equivalence. So such an extension cannot be a benign enrichment of the old scope; it is independent same-domain authority import. These three cases are exhaustive and mutually exclusive.
\end{proof}

\begin{theorem}[Rescue-route exhaustion]
\label{thm:rescue-route-exhaustion}
Every same-domain attempt to rescue a G\"odel-style incompleteness objection by semantic truth, satisfaction, intended-model semantics, reflection, infinitary rule-power, transfinite closure, meta-foundational re-anchoring, or richer same-domain extension falls into exactly one of the following classes:
\begin{enumerate}
  \item closure-inert bookkeeping;
  \item extensional gate renaming;
  \item finite same-domain witness already captured by standing classification or coherent trajectory structure;
  \item explicit scope change; or
  \item inadmissible same-domain authority import.
\end{enumerate}
None of these classes yields a new same-domain governance-level closure threat on the fixed scope.
\end{theorem}

\begin{proof}
Theorem~\ref{thm:no-semantic-standing-import} classifies semantic and meta-foundational routes as inert, gate renaming, scope-changing, or inadmissible. Theorem~\ref{thm:infinitary-admissibility-collapse} classifies infinitary and reflection-based routes: any governance-relevant same-domain force collapses to a finite witness already captured by standing or trajectory structure, and any residual claim without such a witness is inert, scope-changing, or inadmissible. Theorem~\ref{thm:same-domain-extension-trichotomy} classifies richer same-domain extensions as bookkeeping, explicit scope change, or inadmissible authority import. These families exhaust the rescue space. None of the resulting classes yields a new same-domain threat: inert bookkeeping changes no inferential fact, gate renaming adds no new authority, finite witness collapse adds nothing beyond the standing/trajectory structure already fixed, scope change is not same-domain, and inadmissible authority import is excluded by the regime.
\end{proof}


\begin{theorem}[G\"odel threat alignment]
\label{thm:godel-threat-alignment}
Let $F$ be any formal system or coding architecture associated with the fixed domain. If a G\"odel-style undecidability or unprovability phenomenon in $F$ is claimed to threaten governance-level closure on that domain, then that phenomenon must be realized as a same-domain governance-relevant distinction. In particular, it must survive into the standing quotient or coherent trajectory structure of the fixed domain.
\end{theorem}

\begin{proof}
This is an immediate application of Theorem~\ref{thm:inferential-necessity-closure-distinctions}. If the alleged phenomenon does not survive into standing classification or coherent trajectory structure, then it is closure-inert and does not threaten governance-level closure on the fixed domain. Therefore any genuine same-domain G\"odel threat must survive there.
\end{proof}

\begin{theorem}[Translation filter for ordinary G\"odel objections]
\label{thm:translation-filter-ordinary-godel}
Let $F$ be an ordinary recursively axiomatized formal system with an undecidable or unprovable sentence $G_F$. To threaten governance-level closure on the fixed domain, the phenomenon associated with $G_F$ must be translated into a same-domain distinction affecting admissibility, standing, or licensed continuation on that domain. If no such translation exists, then the phenomenon is closure-inert relative to the fixed domain. If such a translation does exist, then it requires the same-domain prerequisite package of Theorem~\ref{thm:same-domain-godel-prerequisites}; any code-based negative fixed-point translation is blocked by Theorems~\ref{thm:no-independent-code-predicate} and \ref{thm:no-same-domain-diagonalization}, and any semantic, infinitary, or extension rescue translation is blocked by Theorem~\ref{thm:rescue-route-exhaustion}. Hence ordinary G\"odel phenomena threaten closure only by a same-domain translation that the present regime excludes.
\end{theorem}

\begin{proof}
If the phenomenon associated with $G_F$ is never translated into a same-domain distinction affecting admissibility, standing, or licensed continuation on the fixed domain, then by Theorem~\ref{thm:godel-threat-alignment} it is closure-inert relative to that domain.

Assume instead that such a same-domain translation exists. Then by Theorem~\ref{thm:godel-threat-alignment} it must survive into standing classification or coherent trajectory structure, and by Theorem~\ref{thm:same-domain-godel-prerequisites} it must realize the same-domain G\"odel prerequisite package. If the translation is code-based and negative-fixed-point in form, then Theorems~\ref{thm:no-independent-code-predicate} and \ref{thm:no-same-domain-diagonalization} eliminate the only remaining governance-relevant self-reference candidate after code-role exhaustion. If the translation instead attempts to rescue the objection by semantic truth, satisfaction, reflection, infinitary rule-power, or richer same-domain extension, then Theorem~\ref{thm:rescue-route-exhaustion} blocks those same-domain rescue routes. Therefore no ordinary G\"odel phenomenon threatens closure on the fixed domain except by a same-domain translation that the present regime forbids.
\end{proof}

\begin{remark}
Ordinary PA-style objections are now subsumed by Theorem~\ref{thm:translation-filter-ordinary-godel}. An undecidable sentence such as $G_F$ threatens closure on the fixed domain only if it is translated into a same-domain distinction affecting admissibility, standing, or licensed continuation; if no such translation occurs, the phenomenon is inert, and if it does occur, it requires the same-domain coding and self-reference package already eliminated above.
\end{remark}

\subsection{Terminal route exhaustion and non-instantiability}

\begin{theorem}[Exhaustion of same-domain incompleteness-objection routes]
\label{thm:incompleteness-route-exhaustion}
Let $\Omega$ be any alleged incompleteness objection to governance-level closure on the fixed scope, offered as a same-domain objection. Then exactly one of the following routes must be in play:
\begin{enumerate}
  \item code-level readback of code judgments as governance-relevant constraints on the original acts or trajectories;
  \item semantic or meta-foundational import;
  \item reflection, infinitary rule-power, or transfinite closure;
  \item a second same-scope equivalence, quotient, or standing-bearing refinement beside admissible standing classification;
  \item a richer same-domain extension or new gate; or
  \item explicit scope change.
\end{enumerate}
The first five routes are impossible or closure-inert on the fixed scope; the sixth is not same-domain. Hence no same-domain incompleteness-objection route remains.
\end{theorem}

\begin{proof}
Any genuine threat must survive into standing classification or coherent trajectory structure by Theorem~\ref{thm:closure-relevance-transmission}. To acquire such same-domain force, the objection must proceed by code-level readback, semantic or meta-foundational import, reflection or infinitary rule-power, a second equivalence or parallel quotient, a richer same-domain extension or new gate, or explicit scope change. No seventh same-domain route remains.

Route (1) is exhausted by Theorems~\ref{thm:code-role-exhaustion}, \ref{thm:no-independent-code-predicate}, and \ref{thm:no-same-domain-diagonalization}. Route (2) is eliminated by Theorem~\ref{thm:no-semantic-standing-import}. Route (3) is eliminated by Theorem~\ref{thm:infinitary-admissibility-collapse}. Route (4) is eliminated by the no-second-same-scope-equivalence theorem of Section~\ref{sec:standing-regime}. Route (5) is eliminated by Theorem~\ref{thm:same-domain-extension-trichotomy} together with the unique same-scope gate. Route (6) is, by definition, not same-domain. Therefore the first five routes are impossible or closure-inert and the sixth falls outside the claim.
\end{proof}

\begin{theorem}[Non-instantiability of G\"odel incompleteness in the unique admissible interior]
\label{thm:godel-non-instantiability}
Under the fixed-domain theorem package already proved in Sections~\ref{sec:modeling-necessity}--\ref{sec:physics-closure}, no same-domain G\"odel threat is instantiable in the unique admissible interior. Equivalently, classical incompleteness remains valid in its standard setting, but a same-domain G\"odel architecture relevant to governance-level closure cannot be instantiated inside the admissibility-governed regime of this capstone.
\end{theorem}

\begin{proof}
Assume, for contradiction, that a same-domain G\"odel threat is instantiable on the fixed scope. By Theorems~\ref{thm:godel-threat-alignment} and \ref{thm:same-domain-godel-prerequisites}, the alleged threat must be realized as a same-domain governance-relevant distinction surviving into the standing quotient or coherent trajectory structure. But Theorem~\ref{thm:incompleteness-route-exhaustion} already exhausts every same-domain incompleteness-objection route and shows that the only remaining route is explicit scope change, which is not same-domain. That contradicts the assumption that a same-domain G\"odel threat is instantiable on the fixed scope. Therefore no such threat exists inside the unique admissible interior.
\end{proof}

\begin{corollary}[No governance-relevant ``true but unprovable'' split on the fixed scope]
There is no same-domain admissibility-relevant positive-side split between ``admitted and provable'' and ``admitted but merely true or unprovable'' on the fixed scope.
\end{corollary}

\begin{proof}
Such a split would be a standing-bearing distinction among admitted acts. By Theorem~\ref{thm:closure-relevance-transmission}, it would have to survive into standing classification or coherent trajectory structure. That makes it either a same-side refinement, a second same-scope equivalence relation, or a new same-domain gate. Each possibility is excluded by the route-exhaustion theorem and the standing theorem family already fixed earlier. So no such split exists.
\end{proof}

\begin{corollary}[Governance-certified results are not vulnerable to same-domain incompleteness objections]
Any result whose certification remains inside the fixed admissibility-governed regime of this capstone is not vulnerable to a same-domain G\"odel incompleteness objection.
\end{corollary}

\begin{proof}
Immediate from Theorem~\ref{thm:godel-non-instantiability}. The point is structural: once same-domain route exhaustion is complete, no internal incompleteness route remains by which the fixed domain could be reopened from above.
\end{proof}

This completes the exhaustion at the level of same-domain counterexamples: no same-domain construction and no regime-exiting construction yields a counterexample to closure.

\paragraph{Kernel-level exhaustion pressure.} The route-exhaustion result of this section should be read together with Theorem~
ef{thm:kernel-exhaustion-apparent-counterexamples}. An apparent incompleteness-based escape either collapses back into the same kernel-governed regime, fails as a lower-generator construction, or exits the target phenomenon altogether. Accordingly the meta-closure layer does not merely answer one objection among others; it closes the remaining same-domain objection space after kernel-level exhaustion has already ruled out rival realizations of the same target.
